\documentclass[10pt,a4paper]{amsart}
\usepackage[margin=19mm]{geometry}
\usepackage{amsmath,amssymb,mathtools}
\usepackage[scr=rsfs,cal=euler]{mathalfa}
\usepackage{xfrac}
\usepackage[unicode,hidelinks,bookmarks=false]{hyperref}
\newcommand{\ie}{i.e.\ }
\newcommand{\cf}{cf.\ }
\newcommand{\resp}{respectively\ }
\newcommand{\Subsection}{\subsection}
\providecommand{\nobreakdash}{\nobreak\hbox{-}}
\setlength{\emergencystretch}{2em}
\multlinegap0pt
\usepackage{amssymb}
\usepackage{tikz}
\newtheorem{thm}{Theorem}
\DeclareMathOperator{\Sing}{Sing}
\def\cC{\mathcal{C}}
\def\cR{{\mathcal R}}
\def\kk{\Bbbk}
\def\w{\omega}
\title{On modular inequalities for plane projective curves}
\author{Jose Ignacio Cogolludo-Agust\'in$^*$, Anca~M\u acinic$^{**}$}
\date{Source: arXiv:2605.26237v1 (CC BY 4.0)}
\begin{document}
\maketitle

\noindent\emph{Source and licence.} On modular inequalities for plane projective curves. Version arXiv:2605.26237v1 (CC BY 4.0), distributed by arXiv under the Creative Commons Attribution 4.0 International licence, http://creativecommons.org/licenses/by/4.0/, as recorded in the arXiv licence metadata for this exact version and shown on the abstract page https://arxiv.org/abs/2605.26237v1. Full licence notice: \texttt{licenses/arxiv-2605.26237-CC-BY-4.0.txt}.\par\smallskip

\noindent\textit{Editorial scope.} Counted unit: the article's thm thm:fibered (source0.tex lines 959--977). The article's complete original proof of it is delivered with it (source0.tex lines 979--1056, one \texttt{proof} environment of 78 lines). No mathematics is written, repaired or re-derived: the statement and the proof are the article's own lines, in the article's own notation, and no retained line is substituted or re-flowed.  Retained uncounted as the source-local context the proof needs: lines 475--478; lines 482--486; lines 492--496; lines 953--956.  Nothing was omitted from any retained span, so the package carries no omission record.  No retained line contains a citation, so the wrapper carries no bibliography and no bibliography entry.  No retained line contains a figure environment, an image-inclusion command or drawing code, so no image is delivered and the package declares no asset dependency.  Editorial formatting only, in addition: an AMS \texttt{amsart} preamble that restates the article's own declarations and macros used by the retained lines, namely the environments \texttt{thm} on separate counters, together with the standard packages those lines need.  The retained environments print in this document as Theorem 1 in order of appearance, whereas the article prints the same items with the numbering of its own full document.  The complete native source of the article as distributed in the arXiv e-print for this exact version is bundled unchanged as \texttt{sources/201404/source0.tex}.
\begin{equation}
\label{eq:awb}
\alpha\wedge \omega=\sum_P \lambda_{\delta,P} \bar\psi_P^\delta + \sum_{i=1}^r \lambda_i \bar\psi^i_\infty.
\end{equation}

\begin{equation}
\label{eq:lambda_i}
    \lambda_i=\sum_{k=1}^r d_k(\alpha_k\beta_i-\alpha_i\beta_k)=
\left| \array{cc} \beta_i & \alpha_i \\ \sum_k d_k\beta_k & \sum_k d_k\alpha_k \endarray\right|.
\end{equation}

\begin{equation}
\label{eq:det}
\left| \array{cc} \beta_i & \alpha_i \\
\sum_k \mu_P(\delta,\cC_k)\beta_k & \sum_k \mu_P(\delta,\cC_k)\alpha_k \endarray \right|=0.
\end{equation}

\begin{equation}
\label{eq:R_phi}
\dim_\kk \cR_{\varphi,\kk} = r.
\end{equation}

\begin{thm}
\label{thm:fibered}
Consider $\cC_\varphi=\cC_1\cup\cdots\cup\cC_r$ a quasi fiber--type curve, $\mathrm{char}(\kk)=p$, 
such that $m_i=0\in\kk$ for all $i=1,...,r$. Assume $\omega:=\sum_{i,j} m_{i,j}\alpha_{i}\sigma_{i,j}\in\cR_{\varphi,\kk}$ is associated with $\cC_\varphi$. Then, 
\begin{enumerate}
    \item\label{thm:fibered:1} If $\sum_i \alpha_i=0\in\kk$,
then $\omega\in \cR_{r-2}(W_\cC,\kk)$, that is,
\begin{equation}
\label{eq:fibered:general}
h^1(A^*_{W_{\cC},\kk},\cdot\wedge\w)\geq r-2 
\end{equation}
    \item\label{thm:fibered:2} If $\cC_\varphi$ has index $k=0\in\kk$,
then $\cR_{\varphi,\kk}\subset \cR_{r-1}(W_\cC,\kk)$, that is,
\begin{equation}
\label{eq:fibered:indexp}
h^1(A^*_{W_{\cC},\kk},\cdot\wedge\w)\geq r-1.
\end{equation}
\end{enumerate}
\end{thm}

\begin{proof}
%We will prove the 
%second part first and will indicate how the first part can be shown. 
We use the notation $G_1, \ldots ,G_r,G^k$ for the polynomials in the pencil given by $\cC_\varphi$ as described at the beginning of this section. Let $\cC_{i,j}: F_{i,j}=0$. 
Given a point $P\in \cC$ and a branch $\delta$ of $\cC_i$, we denote by $\mu_P(\delta,G_i)$ the multiplicity of the intersection of $\delta$ with $G_i$ at $P$. Note that 
$\mu_P(\delta,G^k)=k\mu_P(\delta,G)$.

The singular locus $\Sing(\cC)$ can be decomposed into two disjoint sets $\Sing(\cC)=B\cup N$ according to whether or not they are base points. In particular, we have the following.
\begin{enumerate}
\item \label{item:base}
If $P\in B$ (base points), then $G_i(P)=0$ for all $i \in \{1, \ldots ,r\}$. Moreover, given a branch of $\cC_i$ at $P$ (i.e. $\phi_P(\delta)=i$), then 
\begin{equation} \label{eq:mu_delta}
\mu_P(\delta,G_j)=
\mu_P(\delta,F_j)+\mu_P(\delta,H_j^{m_j})=
\sum_k m_{j,k}\mu_P(\delta,\cC_{j,k})+m_j\mu_P(\delta,H_j)=\mu_\delta
\end{equation}
has the same value regardless of $j\neq i$. Moreover, since $G^k$ is a fiber in the linear system, one has $\mu_\delta=k\mu_P(\delta,G)=0\in\kk$. Since $m_j=0\in\kk$ this implies 
\begin{equation}\label{eq:quasift}
\sum_k m_{j,k}\mu_P(\delta,\cC_{j,k})=0\in\kk \textrm{ for any } j\neq i.
\end{equation}
\item \label{item:nonbase}
If $P\in N$ (non-base points) is such that $\phi_P(\Delta_P)=\{i\}$, then $F_j(P)=0$ if and only if $i=j$.
\end{enumerate}
We will use Conditions~\eqref{eq:lambda_i} and \eqref{eq:det} 
to show that $\w_1\curlywedge \w_2=0$ for any two 1-forms $\w_1,\w_2\in \cR_{\varphi,\kk}$ associated with $\cC_\varphi$, say $\w_1 = \sum_{i,j} m_{i,j}\alpha_{i}\sigma_{i,j}, \; \w_2 = \sum_{i,j} m_{i,j}\beta_{i}\sigma_{i,j}$. 
As in~\eqref{eq:awb} we use the notation 
$$
\w_1\curlywedge\w_2=
\sum_{\delta,P}\lambda_{\delta,P} \bar\psi_{P}^\delta+
\sum_{i,j} \lambda_{i,j}\bar\psi_\infty^{i,j}.
$$
We will show $\lambda_{\delta,P}=\lambda_{i,j}=0$ for all $\delta\in\Delta_P$, all $P\in S$, and $i=1,...,r$.
Consider $P\in N$, then 
condition~\eqref{eq:det} becomes
\begin{equation} \label{eq:non_base_P}
\lambda_{\delta,P}=    
\left|
\array{cc}
\alpha_{i,j} & \beta_{i,j}\\
\sum_s \mu_P(\delta,\cC_{i,s})m_{i,s}\alpha_{i,s} & \sum_s \mu_P(\delta,\cC_{i,s})m_{i,s}\beta_{i,s}\\
\endarray
\right|=0.
\end{equation}
since $\alpha_{i,j}=\alpha_{i,s}=\alpha_i$
and 
$\beta_{i,j}=\beta_{i,s}=\beta_i$ by hypothesis.

Assume now that $P\in B$. In that case, condition~\eqref{eq:det} is also satisfied 
since 
\begin{equation} \label{eq:q_fibr_mod_p}
    \sum_{i,k} \mu_P(\delta,\cC_{i,k})m_{i,k}\alpha_{i,k}=
\mu_\delta \sum_i \alpha_i = 0\in\kk.
\end{equation}
The equalities follow from conditions \eqref{eq:mu_delta} and \eqref{eq:quasift} since $P\in B$. 

Condition~\eqref{eq:lambda_i} becomes
\begin{equation} \label{eq:non_base_P_lambda}
\lambda_i=
\left|
\array{cc}
\alpha_{i,j} & \beta_{i,j}\\
\sum_{s,t} d_{s,t}m_{s,t}\alpha_{s,t} & \sum_{s,t} d_{s,t}m_{s,t}\beta_{s,t}\\
\endarray
\right|=
    \left|
\array{cc}
\alpha_{i} & \beta_{i}\\
kd_G\sum_{s} \alpha_{s} & kd_G\sum_{s} \beta_{s}\\
\endarray
\right|=
0
\end{equation}
since either $k=0\in\kk^*$ or $\sum_{s} \alpha_{s}=\sum_{s} \beta_{s}$ by hypothesis.

The result on the dimension follows by~\eqref{eq:R_phi}.

Note that equality~\eqref{eq:q_fibr_mod_p} is also true if $\sum_i\alpha_i=0\in\kk$ and $m_i=0\in\kk$ which is the hypothesis of the first part. The dimension $r-2$ is given by~\eqref{eq:R_phi} and the extra condition $\sum_i\alpha_i=0\in\kk$.
\end{proof}

\end{document}
